
\textbf{Corollaire 6.3.} \uline{Avec les hypothèses de} 6.1, \uline{soit} $A$ \uline{un anneau commutatif}. \uline{Pour qu'un élément} $y$ \uline{de} $K^{\bullet}(X)\otimes_{\mathbb Z}A=K^{\bullet}(X)_A$ \uline{soit inversible}, \uline{il faut et il suffit que} $z=(\varepsilon\otimes\Id_A)(y)$ \uline{le soit}. \uline{En particulier}, \uline{si} $\underline{E}$ \uline{est un faisceau localement libre de type fini sur} $X$, \uline{partout} $\ne0$, \uline{alors sa classe dans} $K^{\bullet}(X)_{\mathbb Q}$ \uline{est inversible}.

L'idéal d'augmentation de $K^{\bullet}(X)\otimes_{\mathbb Z}A$ est $\Fil^1K^{\bullet}(X)\otimes_{\mathbb Z}A$, et est donc un nilidéal, d'après 6.1. Pour qu'un élément de $K^{\bullet}(X)\otimes_{\mathbb Z}A$ soit inversible il est donc nécessaire et suffisant que son augmentation le soit.

\textbf{Corollaire 6.4.} \uline{Soient} $X$ \uline{un schéma quasi}-\uline{compact}, \uline{ayant un faisceau ample}, \uline{et} $f:X'\to X$ \uline{un morphisme propre et parfait}. \uline{On suppose de plus vérifiée l'une des deux conditions suivantes}:
\[
  \text{i) } f\text{ \uline{projectif}},\qquad
  \text{ii) } X\text{ \uline{noethérien}},
\]
\uline{de façon que le} «\uline{théorème de finitude}» \uline{soit vrai} (III 2.2). \uline{On se donne un anneau} $A$, \uline{et} \uline{on considère l'homomorphisme} «\uline{image directe}» (IV 2.11):
\[
  f_*:K^{\bullet}(X')\otimes_{\mathbb Z}A\longrightarrow K^{\bullet}(X)\otimes_{\mathbb Z}A.
\]
\uline{On suppose enfin qu'il existe un élément} $x'\in K^{\bullet}(X')\otimes_{\mathbb Z}A$ \uline{tel que} $\varepsilon(f_*(x'))\in H^0(X,\underline{A})$ \uline{soit inversible}. \uline{Alors l'homomorphisme} :
\[
  f^*:K^{\bullet}(X)\otimes_{\mathbb Z}A\longrightarrow K^{\bullet}(X')\otimes_{\mathbb Z}A
\]
\uline{est injectif}.

On écrit la formule de projection:
\[
  f_*(f^*(x)x')=x.f_*(x').
\]
Or $f_*(x')$ est d'augmentation inversible, donc est inversible (6.3). Par suite, l'homomorphisme composé :
